\documentclass[aspectratio=169,10pt]{beamer}

% Shared Beamer setup for Fall 2026 course slides.
% Clean Cal Poly-inspired theme: no custom class files or uncommon packages.
\mode<presentation>{
  \usetheme{default}
  \usefonttheme{professionalfonts}
}

\definecolor{CalPolyGreen}{HTML}{154734}
\definecolor{CalPolyGold}{HTML}{BD8B13}
\definecolor{CalPolySage}{HTML}{7FA28F}
\definecolor{CalPolyMint}{HTML}{DDF1E7}
\definecolor{CalPolySand}{HTML}{A29F86}
\definecolor{CalPolyBeige}{HTML}{EFEEDE}
\definecolor{DeckInk}{HTML}{1F2933}
\definecolor{DeckMuted}{HTML}{5F6B7A}
\definecolor{DeckSoft}{HTML}{F7F7F5}

\setbeamersize{text margin left=0.55in,text margin right=0.55in}
\setbeamercolor{normal text}{fg=DeckInk,bg=white}
\setbeamercolor{structure}{fg=CalPolyGreen}
\setbeamercolor{title}{fg=white}
\setbeamercolor{frametitle}{fg=CalPolyGreen,bg=white}
\setbeamercolor{block title}{bg=CalPolySage,fg=white}
\setbeamercolor{block body}{bg=CalPolyMint,fg=black}
\setbeamercolor{alerted text}{fg=CalPolyGold}

\setbeamerfont{title}{size=\Huge,series=\bfseries}
\setbeamerfont{frametitle}{size=\Large,series=\bfseries}
\setbeamerfont{block title}{size=\normalsize,series=\bfseries}

\setbeamertemplate{navigation symbols}{}
\setbeamertemplate{itemize item}{\color{CalPolyGold}$\blacktriangleright$}
\setbeamertemplate{itemize subitem}{\color{CalPolyGreen}$\bullet$}
\setbeamertemplate{footline}{
  \hfill{\color{DeckMuted}\scriptsize\insertframenumber/\inserttotalframenumber}
  \hspace{0.18in}\vspace{0.12in}
}
\setbeamertemplate{frametitle}{
  \vspace{0.18in}
  \hspace{0.02in}\insertframetitle\par
  \vspace{0.08in}
  \noindent\makebox[\linewidth][l]{%
    {\color{CalPolyGold}\rule{0.55in}{2pt}}%
    {\color{CalPolyGreen}\rule{\dimexpr\linewidth-0.55in\relax}{2pt}}%
  }
}

\newcommand{\CourseNumber}{Course}
\newcommand{\CourseTitle}{Course Title}
\newcommand{\LectureNumber}{0}
\newcommand{\LectureTitle}{Lecture Title}
\newcommand{\LectureDate}{Date}
\newcommand{\CourseUnit}{Unit}

\newcommand{\coursenumber}[1]{\renewcommand{\CourseNumber}{#1}}
\newcommand{\coursetitle}[1]{\renewcommand{\CourseTitle}{#1}}
\newcommand{\lecturenumber}[1]{\renewcommand{\LectureNumber}{#1}}
\newcommand{\lecturetitle}[1]{\renewcommand{\LectureTitle}{#1}}
\newcommand{\lecturedate}[1]{\renewcommand{\LectureDate}{#1}}
\newcommand{\courseunit}[1]{\renewcommand{\CourseUnit}{#1}}

\author{Kirk Duran}
\institute{California Polytechnic State University, San Luis Obispo}
\date{\LectureDate}

\newcommand{\makedecktitle}{
  \begingroup
  \setbeamercolor{background canvas}{bg=CalPolyGreen}
  \begin{frame}[plain]
    \vfill
    \begin{center}
      {\usebeamerfont{title}\color{white}\LectureTitle\par}
      \vspace{0.18in}
      {\Large\color{white}Lecture \LectureNumber\par}
    \end{center}
    \vfill
    \noindent{\color{white}\rule{\linewidth}{1.2pt}}\par
    \vspace{0.08in}
    \noindent\colorbox{CalPolyGold}{\makebox[0.24\linewidth][c]{\color{white}\Large\CourseNumber}}
    \hspace{0.12in}{\color{white}\Large Kirk Duran}
  \end{frame}
  \endgroup
}

\newenvironment{keyidea}{\begin{block}{Key idea}}{\end{block}}
\newenvironment{activity}{\begin{block}{In class}}{\end{block}}
\newenvironment{cpdefinition}{\begin{block}{Definition}}{\end{block}}
\newenvironment{exampleblockcp}{
  \begingroup
  \setbeamercolor{block title}{bg=CalPolySand,fg=white}
  \setbeamercolor{block body}{bg=CalPolyBeige,fg=black}
  \begin{block}{Example}
}{
  \end{block}
  \endgroup
}

\newcommand{\imageplaceholder}[2][0.3in]{%
  \noindent\makebox[\linewidth][c]{%
    \fcolorbox{CalPolySage}{CalPolyBeige}{%
      \parbox[c][#1][c]{0.86\linewidth}{%
        \centering
        {\color{DeckMuted}\scriptsize Image placeholder: }%
        {\color{DeckInk}\scriptsize\ttfamily\detokenize{#2}}%
      }%
    }%
  }%
}

\newcommand{\sectionframe}[1]{
  \begingroup
  \setbeamercolor{background canvas}{bg=CalPolyGreen}
  \begin{frame}[plain]
    \vfill
    {\color{white}\LARGE\bfseries #1\par}
    \vspace{0.18in}
    {\color{CalPolyGold}\rule{0.22\paperwidth}{3pt}}
    {\color{white}\rule{0.62\paperwidth}{3pt}}
    \vfill
  \end{frame}
  \endgroup
}

\newcommand{\placeholdernote}{
  \begin{block}{Placeholder}
    Replace these bullets with the final lecture material, examples, and in-class activities.
  \end{block}
}


\usepackage[T1]{fontenc}
\usepackage{lmodern}
\usepackage{amsmath,amssymb}
\usepackage{graphicx}
\pdfstringdefDisableCommands{\def\translate#1{#1}}
\graphicspath{{../assets/lecture-31/}}

% Cal Poly CS 4880, Lecture 31, November 9, 2026.
% Continuation of Lecture 30's style, architecture-first teaching, and transparent results.
% About 50 minutes: motivation 5, clustering 18, representations/PCA 20,
% novelty/SVM 4, synthesis 3. Code is for scikit-learn 1.8.0 (no external downloads).
% Walk through code at a high level; full runnable examples live in the companion script.
% All displayed numeric results and images are reproduced by the companion script
% with random_state=4880 and are NOT predictions of results in another environment.
% Hidden ground-truth labels are used only to COLOR or EVALUATE certain toy figures,
% never to train the unsupervised models.

\newcommand{\lecturevisual}[2][2.15in]{%
  \par\vspace{2pt}{\centering\includegraphics[width=.98\linewidth,height=#1,keepaspectratio]{#2}\par}\vspace{2pt}%
}
\newcommand{\architecturestage}[1]{%
  {\centering\setlength{\fboxsep}{5pt}%
  \fcolorbox{CalPolyGreen}{DeckSoft}{\parbox{.86\linewidth}{\centering\small #1}}\par}%
}
\newcommand{\conceptlead}[1]{{\small\bfseries\color{CalPolyGreen}#1\par}}
\newenvironment{calculation}[1][Mathematical description]{%
  \par\begingroup\setbeamercolor{calculation body}{bg=DeckSoft,fg=DeckInk}%
  \begin{beamercolorbox}[sep=3pt,wd=\linewidth]{calculation body}%
  \footnotesize{\bfseries\color{DeckMuted}#1\par}%
  \setlength{\jot}{0pt}\setlength{\abovedisplayskip}{3pt}\setlength{\belowdisplayskip}{3pt}%
}{\end{beamercolorbox}\endgroup}

\setbeamersize{text margin left=.42in,text margin right=.42in}
\coursenumber{CS 4880}
\coursetitle{Artificial Intelligence}
\lecturenumber{31}
\lecturetitle{Unsupervised Learning: Structure, Clustering, and Representations}
\lecturedate{November 9, 2026}
\courseunit{Learning Agents and Machine Learning}

\begin{document}
\makedecktitle

\section{What Can We Learn Without Labels?}

\begin{frame}{A Dataset Can Be Valuable Without Ground-Truth Answers}
\small
\conceptlead{In supervised learning, labels define a target; in unsupervised learning, structure in the inputs is the target.}
\begin{columns}[T,onlytextwidth]
\begin{column}{.56\textwidth}
\lecturevisual[2.35in]{unlabeled_vs_labels.png}
\end{column}
\begin{column}{.41\textwidth}
\begin{itemize}
  \item We observe feature vectors $x_i\in\mathbb R^d$ without provided class labels $y_i$.
  \item We can ask which observations resemble each other, what varies, and which are unusual.
  \item The colors at right are \emph{hidden synthetic truth}, shown after the fact for teaching.
\end{itemize}
\end{column}
\end{columns}
\begin{keyidea}Discovering a pattern is not the same as discovering its real-world meaning.\end{keyidea}
\end{frame}

\begin{frame}{Three Unsupervised Learning Problems}
\small
\begin{columns}[T,onlytextwidth]
\begin{column}{.32\textwidth}
\begin{block}{Clustering}
Which observations belong to similar groups?\\[5pt]
\textbf{Examples:} K-means, DBSCAN, agglomerative clustering.
\end{block}
\end{column}
\begin{column}{.32\textwidth}
\begin{block}{Representations}
Can we describe an observation with fewer, more useful coordinates?\\[5pt]
\textbf{Examples:} PCA, NMF, t-SNE.
\end{block}
\end{column}
\begin{column}{.32\textwidth}
\begin{block}{Novelty detection}
Does a new observation resemble the training population?\\[5pt]
\textbf{Example:} One-Class SVM.
\end{block}
\end{column}
\end{columns}
\begin{keyidea}We still need a notion of similarity, an inductive assumption, and an honest evaluation.\end{keyidea}
\end{frame}

\begin{frame}{Even Without Labels, We Make Assumptions}
\small
\conceptlead{Similarity depends on the representation and the units we choose.}
\lecturevisual[1.08in]{why_scaling_matters.png}
\begin{columns}[T,onlytextwidth]
\begin{column}{.64\textwidth}
\begin{itemize}
\item Raw Euclidean distance gives equal weight to a unit of temperature and a unit of distance.
\item Standardize nonconstant training features to mean 0 and standard deviation 1; reuse those statistics for new inputs.
\item Rescaling can change neighborhoods, centroids, and clusters.
\end{itemize}
\end{column}
\begin{column}{.33\textwidth}
\begin{calculation}[Euclidean distance]
\[d(x_i,x_j)=\sqrt{\sum_{r=1}^{d}(x_{ir}-x_{jr})^2}.\]
Here $d$ is the number of features; $r$ indexes one coordinate.
\end{calculation}
\end{column}
\end{columns}
\end{frame}

\begin{frame}{How Can We Evaluate an Unsupervised Result?}
\small
\begin{itemize}
\item \textbf{Internal:} is a geometric objective small, or are clusters compact and separated? No class labels required.
\item \textbf{External:} compare with labels reserved for analysis. Adjusted Rand index (ARI) measures agreement regardless of cluster IDs: 1 means identical partitions, around 0 means chance agreement.
\item \textbf{Stability:} do reasonable changes to sampling, initialization, or hyperparameters change the answer?
\item \textbf{Utility:} does the discovered representation help a later task or scientific investigation?
\end{itemize}
\begin{keyidea}A compelling two-dimensional plot is neither proof of correctness nor evidence of causation.\end{keyidea}
\end{frame}

\section{Clustering: K-Means and Alternatives}

\begin{frame}{K-Means: Build Groups Around Representative Centers}
\small
\conceptlead{K-means assumes a useful grouping can be described with $k$ centroids.}
\begin{enumerate}
\item Choose the number of groups $k$ and initialize $k$ centers.
\item \textbf{Assignment:} send each point to its nearest center.
\item \textbf{Update:} move each center to the mean of its assigned points.
\item Repeat until assignments or the objective stabilize.
\end{enumerate}
\begin{calculation}[K-means objective, also called inertia]
\[J=\sum_{i=1}^{n}\|x_i-\mu_{z_i}\|_2^2.\]
$z_i\in\{1,\ldots,k\}$ is the cluster assignment of point $i$; $\mu_j$ is centroid $j$.
The algorithm reduces $J$ iteratively but need not find its global minimum.
\end{calculation}
\end{frame}

\begin{frame}{K-Means in Motion: Assign, Update, Repeat}
\small
\conceptlead{The optimization alternates between assigning points and moving the centers.}
\lecturevisual[2.55in]{kmeans_iteration.png}
\begin{itemize}
\item The large \textsf{X} markers denote the current centroid locations.
\item Initialization matters; scikit-learn can run several starts and keep the best solution.
\item A cluster ID is only a numerical index, not a semantic category such as ``safe.''
\end{itemize}
\end{frame}

\begin{frame}[fragile]{Demo 1: Fit K-Means with scikit-learn}
\small
\conceptlead{The labels are generated as a teaching check, but never passed to K-means.}
{\footnotesize\begin{verbatim}
from sklearn.datasets import make_blobs
from sklearn.cluster import KMeans
from sklearn.metrics import silhouette_score, adjusted_rand_score

X, hidden_y = make_blobs(
    n_samples=450,
    centers=[[-5,-1], [0,3], [5,-.7]],
    cluster_std=[.82, 1.17, .9], random_state=4880)
model = KMeans(n_clusters=3, n_init=10, random_state=4880)
cluster_id = model.fit_predict(X)
print(round(model.inertia_, 1))
print(round(silhouette_score(X, cluster_id), 3))
print(round(adjusted_rand_score(hidden_y, cluster_id), 3))
\end{verbatim}}
\begin{keyidea}K-means sees only $X$; adjusted Rand index (ARI) is a post-hoc teaching diagnostic.\end{keyidea}
\end{frame}

\begin{frame}{Demo 1 Results: Three Groups in 450 Points}
\small
\begin{columns}[T,onlytextwidth]
\begin{column}{.58\textwidth}
\lecturevisual[2.55in]{kmeans_on_blobs.png}
\end{column}
\begin{column}{.39\textwidth}
\begin{block}{Measured run}
\textbf{Inertia:} 897.3\\[5pt]
\textbf{Silhouette:} 0.705\\[5pt]
\textbf{Post-hoc ARI:} 0.987\\[5pt]
\textbf{Cluster sizes:} 151, 151, 148
\end{block}
\small \textbf{Interpretation:} the compact, well-separated synthetic groups match K-means' geometric assumptions.
\end{column}
\end{columns}
\end{frame}

\begin{frame}{How Many Clusters? Why Inertia Alone Is Not Enough}
\small
\conceptlead{The algorithm requires $k$, but the dataset does not provide a uniquely correct choice.}
\begin{itemize}
\item More clusters almost always reduce inertia: eventually each point could be its own center.
\item The \textbf{elbow} heuristic looks for diminishing improvements in inertia.
\item \textbf{Silhouette} compares a point's average distance within its cluster with distance to the nearest other cluster.
\item Subject-matter constraints and stability may matter more than maximizing a score.
\end{itemize}
\begin{calculation}[Silhouette of an individual observation]
\[s(i)=\frac{b(i)-a(i)}{\max\{a(i),b(i)\}}\quad\in[-1,1].\]
$a(i)$ is average within-cluster distance; $b(i)$ is the smallest average distance to another cluster. A larger value indicates stronger separation under this distance.
\end{calculation}
\end{frame}

\begin{frame}[fragile]{Demo 2: Try Several Values of $k$}
\small
\conceptlead{Run the same clustering method with different numbers of groups.}
{\footnotesize\begin{verbatim}
from sklearn.cluster import KMeans
from sklearn.metrics import silhouette_score

for k in range(2, 8):
    m = KMeans(n_clusters=k, n_init=10,
               random_state=4880).fit(X)
    score = silhouette_score(X, m.labels_)
    print(k, round(m.inertia_, 1), round(score, 3))
\end{verbatim}}
\begin{calculation}[Printed output: k, inertia, mean silhouette]
\texttt{2: 3740.7, 0.562}\quad \texttt{3: 897.3, 0.705}\quad \texttt{4: 742.6, 0.592}\\[2pt]
\texttt{5: 640.0, 0.578}\quad \texttt{6: 544.6, 0.475}\quad \texttt{7: 483.9, 0.331}
\end{calculation}
\end{frame}

\begin{frame}{Demo 2 Results: Inertia Versus Silhouette}
\small
\lecturevisual[2.7in]{k_selection.png}
\begin{keyidea}For this dataset, the mean silhouette peaks at $k=3$. No universal scoring rule recovers the ``true'' number of real-world groups.\end{keyidea}
\end{frame}

\begin{frame}{The Geometry of the Data Can Defeat K-Means}
\small
\conceptlead{Not all meaningful clusters look like compact blobs around a center.}
\lecturevisual[2.55in]{moons_algorithms.png}
\begin{itemize}
\item These interlocking crescents have two connected structures, not two spherical clouds.
\item A centroid-based partition cuts across the curves; density connections can follow them.
\end{itemize}
\end{frame}

\begin{frame}{DBSCAN: Define a Group by Density, Not a Center}
\small
\conceptlead{Density-Based Spatial Clustering of Applications with Noise identifies dense, connected regions.}
\begin{itemize}
\item \textbf{$\varepsilon$ (eps):} the neighborhood radius around a point.
\item \textbf{\texttt{min\_samples}:} the minimum nearby sample count, including the point itself, for a core point.
\item Core points connect through chains of dense neighborhoods; border points may join nearby groups.
\item Points not assigned to a group are labeled \texttt{-1} (noise).
\item It does not require choosing $k$ in advance, but assumes a useful density scale.
\end{itemize}
\begin{keyidea}DBSCAN can discover nonconvex groups; it can struggle with differing densities or poorly selected radius.\end{keyidea}
\end{frame}

\begin{frame}[fragile]{Demo 3: K-Means Versus DBSCAN on Two Moons}
\small
{\footnotesize\begin{verbatim}
from sklearn.datasets import make_moons
from sklearn.preprocessing import StandardScaler
from sklearn.cluster import KMeans, DBSCAN
from sklearn.metrics import adjusted_rand_score, silhouette_score

X_moon, hidden_y = make_moons(n_samples=440, noise=.07,
                              random_state=4880)
X_moon = StandardScaler().fit_transform(X_moon)
k = KMeans(n_clusters=2, n_init=10, random_state=4880)
d = DBSCAN(eps=.22, min_samples=6)
a, b = k.fit_predict(X_moon), d.fit_predict(X_moon)
print(adjusted_rand_score(hidden_y, a))
print(adjusted_rand_score(hidden_y, b))
print(silhouette_score(X_moon, a))
print(silhouette_score(X_moon, b))
\end{verbatim}}
\begin{calculation}[Rounded output: K-means, then DBSCAN]
\textbf{Post-hoc ARI:} 0.457 vs. 0.986.\qquad
\textbf{Silhouette:} 0.491 vs. 0.289 (noise label included).
\end{calculation}
\end{frame}

\begin{frame}{A Critical Result: The Better Silhouette Is Misleading}
\small
\conceptlead{The internal score prefers K-means, while external knowledge strongly favors DBSCAN.}
\begin{columns}[T,onlytextwidth]
\begin{column}{.57\textwidth}
\lecturevisual[2.5in]{moons_algorithms.png}
\end{column}
\begin{column}{.40\textwidth}
\begin{itemize}
\item \textbf{ARI} compares clusters with synthetic truth only after fitting.
\item \textbf{Silhouette} rewards compact, separated partitions in the selected metric.
\item Neither number tells us what the groups mean or whether they cause an outcome.
\end{itemize}
\end{column}
\end{columns}
\end{frame}

\begin{frame}{DBSCAN Is Sensitive to Its Neighborhood Scale}
\small
\conceptlead{Different density thresholds can split, recover, or merge structure.}
\lecturevisual[2.7in]{dbscan_eps_sensitivity.png}
\begin{itemize}
\item On these standardized two-moon features, \texttt{min\_samples=6} throughout.
\item \texttt{eps=0.12}: 19 groups, 72 noise; \texttt{eps=0.22}: 2 groups, 3 noise; \texttt{eps=0.38}: 2 groups, 0 noise.
\end{itemize}
\end{frame}

\begin{frame}{Agglomerative Clustering: Preserve a Hierarchy}
\small
\conceptlead{Instead of choosing a single partition immediately, repeatedly merge groups.}
\begin{columns}[T,onlytextwidth]
\begin{column}{.62\textwidth}
\lecturevisual[2.6in]{hierarchical_dendrogram.png}
\end{column}
\begin{column}{.35\textwidth}
\begin{itemize}
\item Begin with individual points.
\item Merge the best pair until one hierarchy remains.
\item Different \emph{linkage} rules define which merge is best.
\item This figure uses Ward linkage on 22 sampled points.
\end{itemize}
\end{column}
\end{columns}
\end{frame}

\begin{frame}[fragile]{A Probabilistic Alternative: Gaussian Mixture Models}
\small
\conceptlead{A mixture model gives each point a degree of membership in every component.}
\begin{columns}[T,onlytextwidth]
\begin{column}{.49\textwidth}
{\scriptsize\begin{verbatim}
from sklearn.mixture import GaussianMixture
from sklearn.datasets import make_blobs
X_overlap, _ = make_blobs(
    n_samples=450,
    centers=[[-1.55,-.75],[1.5,-.65],[0,1.6]],
    cluster_std=[1.,1.,1.05],
    random_state=4880)
g = GaussianMixture(n_components=3,
    n_init=5, random_state=4880).fit(X_overlap)
prob = g.predict_proba(X_overlap)
\end{verbatim}}
\end{column}
\begin{column}{.48\textwidth}
\lecturevisual[2.15in]{gmm_soft_membership.png}
\end{column}
\end{columns}
\begin{calculation}[Mixture density]
\[p(x)=\sum_{j=1}^{k}\pi_j\,\mathcal N(x\mid\mu_j,\Sigma_j),\qquad \sum_j\pi_j=1.\]
$\pi_j$ are mixing weights; $\mu_j$ and $\Sigma_j$ describe each Gaussian's center and spread. Memberships are conditional on this fitted model, not guaranteed real-world classes.
\end{calculation}
\end{frame}

\begin{frame}{When Should We Trust an Unsupervised Cluster?}
\small
\begin{itemize}
\item \textbf{State the intended similarity:} Euclidean distance, cosine similarity, connectivity, or a probabilistic model?
\item \textbf{Inspect scale and preprocessing:} one large-unit feature can dominate distance.
\item \textbf{Check alternative models:} K-means, density-based, hierarchical, and mixture assumptions differ.
\item \textbf{Inspect stability:} resampling and changing seeds or parameters can reveal brittle conclusions.
\item \textbf{Validate meaning:} a group of deliveries is not automatically a class called ``dangerous.''
\end{itemize}
\begin{keyidea}Cluster assignments are hypotheses about structure, not labels supplied by nature.\end{keyidea}
\end{frame}

\section{Dimensionality Reduction and Principal Components}

\begin{frame}{High Dimensional Does Not Always Mean High Information}
\small
\conceptlead{Many measured features may be redundant or correlated.}
\begin{columns}[T,onlytextwidth]
\begin{column}{.48\textwidth}
\lecturevisual[2.35in]{digits_samples.png}
\end{column}
\begin{column}{.49\textwidth}
\begin{itemize}
\item Each scikit-learn digit is an $8\times8$ image, represented by 64 pixel features.
\item We may seek a smaller vector $z\in\mathbb R^r$, $r<d$, retaining useful variation.
\item A compressed representation can make visualization, storage, or later learning easier.
\item Compression always introduces a tradeoff: what is retained, and what is lost?
\end{itemize}
\end{column}
\end{columns}
\end{frame}

\begin{frame}{PCA: Find Directions Along Which the Data Vary Most}
\small
\conceptlead{Principal Component Analysis chooses new orthogonal axes, ordered by captured variance.}
\begin{columns}[T,onlytextwidth]
\begin{column}{.57\textwidth}
\lecturevisual[2.2in]{pca_geometry.png}
\end{column}
\begin{column}{.40\textwidth}
\begin{itemize}
\item PC1 is the direction with the largest variance after centering.
\item PC2 is perpendicular to PC1 and explains the largest remaining variance.
\item Retaining fewer axes projects data into fewer dimensions.
\item The axes are computed without outcome labels.
\end{itemize}
\end{column}
\end{columns}
\end{frame}

\begin{frame}{PCA: Center, Project, and Reconstruct}
\small
\conceptlead{Subtract the feature means, keep the strongest directions of variation, and map back to approximate the input.}
\begin{calculation}[First principal direction]
\[v_1=\underset{\|v\|_2=1}{\arg\max}\;\operatorname{Var}(X_c v),\qquad X_c=X-\mathbf{1}\mu^\top.\]
$X\in\mathbb R^{n\times d}$ holds $n$ examples and $d$ features; $\mu$ is the vector of column means; $v_1$ is a unit vector. Subsequent directions are orthogonal to earlier ones.
\end{calculation}
\begin{calculation}[Project and approximately reconstruct]
\[Z=X_cV_r,\qquad \widehat X=ZV_r^\top+\mathbf{1}\mu^\top.\]
$V_r\in\mathbb R^{d\times r}$ contains the first $r$ principal directions; $Z\in\mathbb R^{n\times r}$ is the lower-dimensional representation.
\end{calculation}
\begin{itemize}
\item scikit-learn computes the directions numerically (SVD or a covariance-based solver); it \emph{centers} but does not automatically standardize input features.
\end{itemize}
\end{frame}

\begin{frame}[fragile]{Demo 4: PCA on Four Measurements of Iris Flowers}
\small
\conceptlead{We fit PCA without species labels; colors in the final plot are added after fitting.}
{\footnotesize\begin{verbatim}
from sklearn.datasets import load_iris
from sklearn.preprocessing import StandardScaler
from sklearn.decomposition import PCA

iris = load_iris()
X = StandardScaler().fit_transform(iris.data)
pca = PCA(n_components=4).fit(X)
Z = pca.transform(X)
print(pca.explained_variance_ratio_.round(4))
print(round(100 * sum(pca.explained_variance_ratio_[:2]), 1))
# Plot Z[:, 0], Z[:, 1]; colors use iris.target only AFTER fit.
\end{verbatim}}
\begin{calculation}[Rounded output]
\texttt{[0.7296, 0.2285, 0.0367, 0.0052]}\qquad First 2 PCs: \textbf{95.8\%} of standardized-feature variance.
\end{calculation}
\end{frame}

\begin{frame}{Demo 4 Results: A Two-Dimensional Iris Representation}
\small
\lecturevisual[2.8in]{iris_pca_projection_variance.png}
\begin{keyidea}Two coordinates preserve 95.8\% of total variance \emph{in the standardized four-feature dataset}. Variance retained is not the same as label accuracy or physical meaning.\end{keyidea}
\end{frame}

\begin{frame}{Explained Variance Helps Choose the Dimension}
\small
\conceptlead{A lower-dimensional representation retains a fraction of the original variability.}
\begin{itemize}
\item \textbf{Explained variance ratio} $\lambda_j / \sum_\ell\lambda_\ell$ measures each component's fraction of total variance.
\item \textbf{Cumulative explained variance} adds those fractions for the retained components.
\item Choosing $r$ based on 90\%, 95\%, or another threshold is a design decision, not a guarantee of task performance.
\item PCA can miss nonlinear geometry or a low-variance feature that is important for a prediction task.
\end{itemize}
\begin{calculation}[Iris, standardized first two components]
\[0.7296+0.2285\approx0.9581.\]
The approximation agrees with the plotted \textbf{95.8\%}. The remaining two components account for about 4.2\%.
\end{calculation}
\end{frame}

\begin{frame}{PCA Also Defines a Reconstruction Experiment}
\small
\conceptlead{To ask what information the representation preserves, reconstruct the original signal.}
\lecturevisual[1.9in]{digits_pca_reconstruction.png}
\begin{itemize}
\item More components improve reconstruction, but do not guarantee human recognizability or classifier accuracy.
\end{itemize}
\end{frame}

\begin{frame}[fragile]{Demo 5: Compress and Reconstruct Handwritten Digits}
\small
{\footnotesize\begin{verbatim}
from sklearn.datasets import load_digits
from sklearn.decomposition import PCA
from sklearn.metrics import mean_squared_error

digits = load_digits()
X = digits.data                  # 1797 x 64, no labels needed
for r in [2, 8, 16, 32]:
    pca = PCA(n_components=r, svd_solver='full').fit(X)
    compressed = pca.transform(X)
    reconstruction = pca.inverse_transform(compressed)
    print(r, round(mean_squared_error(X, reconstruction), 2))
\end{verbatim}}
\begin{calculation}[Measured reconstruction mean squared error, in pixel-value units squared]
\texttt{2 PCs: 13.42}\quad\texttt{8 PCs: 6.12}\quad\texttt{16 PCs: 2.83}\quad\texttt{32 PCs: 0.63}
\end{calculation}
\begin{keyidea}This is \emph{in-sample} reconstruction, not held-out classification or evidence of semantic understanding.\end{keyidea}
\end{frame}

\begin{frame}{Representations Can Be Learned in Other Ways}
\small
\conceptlead{PCA is only one way to define a useful lower-dimensional code.}
\begin{itemize}
\item \textbf{PCA:} signed, orthogonal linear components maximize retained variance.
\item \textbf{Nonnegative Matrix Factorization (NMF):} when input values are nonnegative, learn nonnegative additive basis patterns.
\item \textbf{Autoencoders:} neural networks learn an encoder (input to code) and decoder (code to reconstruction); the reconstruction loss trains both. Lecture 30 introduced nonlinear architectures.
\item \textbf{Task-specific representation:} features may be useful because they help downstream prediction, not simply because they reconstruct well.
\end{itemize}
\begin{calculation}[NMF factorization]
\[X\approx WH,\qquad W\ge0,\ H\ge0.\]
$W\in\mathbb R^{n\times r}$ gives example coefficients; $H\in\mathbb R^{r\times d}$ contains learned additive basis vectors.
\end{calculation}
\end{frame}

\begin{frame}[fragile]{Demo 6: Learn Additive Digit Patterns with NMF}
\small
\begin{columns}[T,onlytextwidth]
\begin{column}{.40\textwidth}
{\scriptsize\begin{verbatim}
from sklearn.decomposition import NMF
from sklearn.datasets import load_digits
import numpy as np
X = load_digits().data   # all >= 0
m = NMF(n_components=10,
        init='nndsvda',
        random_state=4880,
        max_iter=1200)
W = m.fit_transform(X)
H = m.components_
relative_error = (
    np.linalg.norm(X - W @ H)
    / np.linalg.norm(X))
print(round(relative_error, 3))
\end{verbatim}}
\end{column}
\begin{column}{.57\textwidth}
\lecturevisual[2.35in]{nmf_digits_components.png}
\end{column}
\end{columns}
\begin{keyidea}Measured relative Frobenius error: \textbf{0.326}. Components are learned image patterns, not necessarily digit classes or human-interpretable parts.\end{keyidea}
\end{frame}

\begin{frame}{t-SNE: A Nonlinear Visualization of Neighborhoods}
\small
\conceptlead{t-distributed Stochastic Neighbor Embedding seeks a low-dimensional map that preserves local similarities.}
\begin{itemize}
\item PCA maximizes global variance along linear directions.
\item t-SNE instead tries to place near-neighbors in the original space close together in the map.
\item It uses an optimization objective based on probability distributions of pairwise similarities.
\item \textbf{Perplexity} changes the effective neighborhood scale; initialization and random seed can change the picture.
\item It is primarily a visualization method; distances \emph{between} separated map groups can be misleading.
\end{itemize}
\begin{keyidea}Seeing discrete islands in a t-SNE figure does not prove that the data contain those many natural clusters.\end{keyidea}
\end{frame}

\begin{frame}[fragile]{Demo 7: Visualize the Same Digits with PCA and t-SNE}
\small
{\footnotesize\begin{verbatim}
import numpy as np
from sklearn.datasets import load_digits
from sklearn.decomposition import PCA
from sklearn.manifold import TSNE
D = load_digits()
idx = np.random.default_rng(4880).choice(
    len(D.data), size=1150, replace=False)
X = D.data[idx]
X40 = PCA(n_components=40, random_state=4880).fit_transform(X)
Z_pca = X40[:, :2]
Z_tsne = TSNE(n_components=2, perplexity=30,
              learning_rate='auto', init='pca',
              max_iter=650, random_state=4880).fit_transform(X40)
# Plot both; use D.target[idx] ONLY to color points afterward.
\end{verbatim}}
\begin{keyidea}Both embeddings learn from image features alone, not the digit identities.\end{keyidea}
\end{frame}

\begin{frame}{Demo 7 Results: Same Inputs, Different Geometry}
\small
\lecturevisual[2.9in]{digits_pca_vs_tsne.png}
\begin{itemize}
\item Digit labels were used solely to color the points \emph{after} embedding.
\item Global distances and cluster areas in the t-SNE map should not be read as original-space distances or densities.
\end{itemize}
\end{frame}

\begin{frame}{A Good Representation Must Generalize Beyond Its Plot}
\small
\conceptlead{Evaluate a transformation on data it did not help fit.}
\lecturevisual[2.0in]{train_test_pipeline.png}
\begin{itemize}
\item Fit scalers, PCA, and any learned transformation on the \textbf{training split}.
\item Apply the frozen transformation to validation or test observations.
\item Choose representation size using validation data, then retrain on development data.
\item Test once after the choices are fixed; otherwise the test set silently becomes tuning data.
\end{itemize}
\end{frame}

\begin{frame}[fragile]{Demo 8: PCA as an Unsupervised Step in a Supervised Pipeline}
\scriptsize
\conceptlead{Fit PCA without labels; evaluate its features using a supervised classifier.}
\begin{verbatim}
from sklearn.datasets import load_digits
from sklearn.pipeline import make_pipeline
from sklearn.preprocessing import StandardScaler
from sklearn.decomposition import PCA
from sklearn.svm import SVC
from sklearn.model_selection import train_test_split
from sklearn.metrics import accuracy_score
X, y = load_digits(return_X_y=True)
Xdev, Xtest, ydev, ytest = train_test_split(
    X, y, test_size=.20, stratify=y, random_state=4880)
Xtr, Xval, ytr, yval = train_test_split(
    Xdev, ydev, test_size=.25, stratify=ydev, random_state=4880)
def candidate(components):
    stages = [StandardScaler()]
    if components is not None:
        stages += [PCA(n_components=components, random_state=4880)]
    stages += [SVC(C=5, gamma='scale', kernel='rbf')]
    return make_pipeline(*stages)
\end{verbatim}
\end{frame}

\begin{frame}[fragile]{Demo 8: Select on Validation, Then Test Once}
\small
{\footnotesize\begin{verbatim}
choices = [None, 8, 16, 32, 48]   # None => original 64 pixels
validation = {}
for r in choices:
    model = candidate(r).fit(Xtr, ytr)
    validation[r] = accuracy_score(yval, model.predict(Xval))
best_r = max(choices, key=lambda r: validation[r])
final = candidate(best_r).fit(Xdev, ydev)
print(best_r, round(accuracy_score(ytest,
                                    final.predict(Xtest)), 4))
\end{verbatim}}
\begin{calculation}[Measured split: 1077 train, 360 validation, 360 test]
Best validation score is tied by \textbf{raw 64 pixels} and \textbf{48 PCs}, each at \textbf{97.5\%}.
The tie breaks toward the first candidate: \textbf{raw pixels}. After retraining, \textbf{final test accuracy = 97.78\%}.
\end{calculation}
\end{frame}

\begin{frame}{Demo 8 Results: PCA Does Not Always Help Prediction}
\small
\lecturevisual[1.9in]{pca_downstream_svc_accuracy.png}
\begin{itemize}
\item PCA reduces features but can discard information a classifier uses.
\item Validation selects an architecture; the test score is only reported after selection and refitting.
\item We should measure runtime, storage, and stability if compression is the objective.
\end{itemize}
\end{frame}

\section{Novelty Detection and the SVM Distinction}

\begin{frame}{Where Does SVM Fit? Standard SVM Is Supervised}
\small
\conceptlead{An ordinary support vector classifier requires class labels. One-Class SVM learns an inlier region.}
\begin{columns}[T,onlytextwidth]
\begin{column}{.48\textwidth}
\begin{block}{\texttt{sklearn.svm.SVC}}
\textbf{Supervised classification.} Fit pairs $(X,y)$; seek a wide separating margin, with a kernel allowing nonlinear boundaries. Used for digits in Demo 8.
\end{block}
\end{column}
\begin{column}{.48\textwidth}
\begin{block}{\texttt{sklearn.svm.OneClassSVM}}
\textbf{Novelty detection.} Fit examples assumed to represent normal operation; estimate a boundary around their support.
\end{block}
\end{column}
\end{columns}
\begin{itemize}
\item This distinction matters: a One-Class SVM does \emph{not} magically label multiple unknown classes.
\item It tests resemblance under a kernel and scale. Predictions are $+1$ (inlier) or $-1$ (outlier); $\nu$ controls the training-error/support-vector tradeoff, not a guaranteed future anomaly rate.
\end{itemize}
\end{frame}

\begin{frame}[fragile]{Demo 9: Fit a One-Class SVM to Normal Data}
\small
{\footnotesize\begin{verbatim}
import numpy as np
from sklearn.svm import OneClassSVM
rng = np.random.default_rng(4880)
normal_train = rng.normal(0, .6, size=(240, 2))
normal_test = rng.normal(0, .6, size=(80, 2))
anomaly_test = rng.uniform(2, 3.6, size=(50, 2)) * \
               rng.choice([-1, 1], size=(50, 2))
m = OneClassSVM(kernel='rbf', gamma=.4, nu=.05)
m.fit(normal_train)
print((m.predict(normal_test) == 1).mean())
print((m.predict(anomaly_test) == -1).sum())
\end{verbatim}}
\begin{calculation}[Measured output]
\textbf{Normal acceptance:} 0.925 (74 of 80).\qquad
\textbf{Synthetic anomalies flagged:} 50 of 50.
\end{calculation}
\end{frame}

\begin{frame}{Demo 9 Results: Novelty Is Not a Diagnosis}
\small
\begin{columns}[T,onlytextwidth]
\begin{column}{.57\textwidth}
\lecturevisual[2.75in]{one_class_svm_boundary.png}
\end{column}
\begin{column}{.40\textwidth}
\begin{itemize}
\item The green contour is the learned \emph{decision boundary} around training inliers.
\item Outliers were deliberately placed far away; this is an \emph{easy synthetic test}.
\item ``Outside normal'' is not proof of an equipment fault, fraud, or danger.
\item Real systems must measure false alarms under deployment shift.
\end{itemize}
\end{column}
\end{columns}
\end{frame}

\section{Synthesis and Research Practice}

\begin{frame}{Match the Question to the Unsupervised Method}
\footnotesize
\begin{tabular}{@{}p{.22\linewidth}p{.31\linewidth}p{.39\linewidth}@{}}
\textbf{Method} & \textbf{What it looks for} & \textbf{Main assumption / caveat}\\[2pt]\hline
K-means & $k$ centroid-based partitions & Compact groups; choose $k$; sensitive to scaling\\[2pt]
DBSCAN & Connected dense regions & Choose density scale; differing densities are hard\\[2pt]
Agglomerative & Nested merges / hierarchy & Linkage changes the hierarchy\\[2pt]
Gaussian mixture & Soft component memberships & Model family and number of components matter\\[2pt]
PCA & Low-rank linear projection & Maximizes variance, not predictive value\\[2pt]
NMF & Additive nonnegative factors & Requires $X\ge0$; factors need not be semantic\\[2pt]
t-SNE & Low-dimensional neighbor visualization & Plot geometry is not a reliable global metric\\[2pt]
One-Class SVM & Boundary of familiar observations & Sensitive to kernel/scale and training purity
\end{tabular}
\vspace{1pt}
\begin{keyidea}Choose the method for the scientific or operational question, not for the most visually striking figure.\end{keyidea}
\end{frame}

\begin{frame}{A Practical Checklist Before Claiming Discovery}
\small
\begin{enumerate}
\item Define which variables exist and which labels, if any, are missing.
\item Document units, scaling, distance metric, randomness, and model assumptions.
\item Separate fitting from evaluation; do not let future or test data alter preprocessing.
\item Compare plausible alternatives and assess parameter sensitivity.
\item Distinguish geometric scores, post-hoc external labels, and actual task utility.
\item Have a domain expert inspect whether the representation is operationally meaningful.
\end{enumerate}
\begin{keyidea}Unsupervised learning finds structure in observations. Humans must still test what that structure means.\end{keyidea}
\end{frame}

\begin{frame}{Connecting Lecture 31 Back to Learning Agents}
\small
\conceptlead{A delivery robot can collect far more sensor history than people can manually label.}
\begin{itemize}
\item \textbf{Clustering:} organize similar traversal segments for human review.
\item \textbf{Dimensionality reduction:} compress correlated sensor features into a smaller learned state description.
\item \textbf{Anomaly detection:} flag an unusual motor or navigation signature for inspection.
\item \textbf{Supervised follow-up:} once humans label some examples, evaluate whether the discovered features predict real outcomes.
\end{itemize}
\begin{keyidea}Do not treat a cluster ID as a safety label or a novelty flag as a confirmed mechanical failure.\end{keyidea}
\end{frame}

\begin{frame}{References, Implementation, and Reproducibility}
\scriptsize
\begin{itemize}
\item scikit-learn User Guide: \href{https://scikit-learn.org/stable/modules/clustering.html}{Clustering}; \href{https://scikit-learn.org/stable/modules/decomposition.html}{Matrix decomposition and PCA}; \href{https://scikit-learn.org/stable/modules/manifold.html}{Manifold learning}; \href{https://scikit-learn.org/stable/modules/outlier_detection.html}{Novelty and outlier detection}.
\item scikit-learn \href{https://scikit-learn.org/stable/modules/generated/sklearn.svm.OneClassSVM.html}{OneClassSVM} and \href{https://scikit-learn.org/stable/modules/generated/sklearn.pipeline.Pipeline.html}{Pipeline} API documentation.
\item MacQueen (1967), \emph{Some Methods for Classification and Analysis of Multivariate Observations}; Ester et al. (1996), \emph{A Density-Based Algorithm for Discovering Clusters in Large Spatial Databases with Noise}.
\item Pearson (1901), \emph{On Lines and Planes of Closest Fit to Systems of Points in Space}; van der Maaten and Hinton (2008), \emph{Visualizing Data using t-SNE}; Sch\"olkopf et al. (2001), \emph{Estimating the Support of a High-Dimensional Distribution}.
\item Reproducibility: the companion \texttt{generate-figures.py} regenerates all 18 included figures and numerical outputs from built-in scikit-learn datasets and synthetic data; no web access required. Verified using scikit-learn 1.8.0, NumPy 2.3.5, Matplotlib 3.10.8, and SciPy 1.18.1; dependencies are pinned beside the script.
\item All numeric outcomes are \emph{measured on the supplied run}; different library versions or environments can produce slightly different fits or embeddings. Figures with ground-truth coloring never used that coloring to fit an unsupervised model.
\end{itemize}
\end{frame}
\end{document}
